% =============================================================================
% APPENDICE: Bibliografia scientifica ordinata
% Formato: Autore (anno). Titolo. Rivista/Editore. DOI quando stabile.
% =============================================================================

\chapter*{Bibliografia scientifica}
\addcontentsline{toc}{chapter}{Bibliografia scientifica}

Le note a piè di pagina citano tutte le fonti usate nel testo. Questa bibliografia raccoglie, in ordine alfabetico, i lavori scientifici su cui poggiano le affermazioni principali del libro, con il DOI quando è stabile. La precede una tabella che indica, per ciascuna di quelle affermazioni, il grado di evidenza: è il modo più onesto che conosco per distinguere ciò che sappiamo da ciò che crediamo di sapere.

\section*{Il grado di evidenza delle affermazioni principali}

\begin{concetto}[Legenda]
\textbf{Solida}: risultato replicato in molti studi indipendenti, spesso con meta-analisi.\\
\textbf{Buona}: risultati coerenti, ma con limiti di campione, contesto o metodo.\\
\textbf{Discussa}: risultati contrastanti, repliche fallite o forte dibattito tra esperti.\\
\textbf{Ipotesi}: idea plausibile, oppure metafora o scenario, senza prove dirette sufficienti.
\end{concetto}

{\small
\renewcommand{\arraystretch}{1.25}
\begin{longtable}{p{0.62\textwidth}p{0.12\textwidth}p{0.14\textwidth}}
\toprule
\textbf{Affermazione} & \textbf{Grado} & \textbf{Capitolo}\\
\midrule
\endhead
La memoria ricostruisce i ricordi e può incorporare informazioni successive & Solida & \ref{cap:p1c6}\\
Il bias di conferma è diffuso in contesti molto diversi & Solida & \ref{cap:p1c7}\\
Vediamo i bias negli altri più che in noi (punto cieco) & Solida & \ref{cap:p1c7}, \ref{cap:p4c25}\\
Ancoraggio e negligenza del tasso di base influenzano stime e giudizi & Solida & \ref{cap:p2c14}\\
Il «cervello trino» descrive l'anatomia evolutiva del cervello & Superata (solo metafora) & \ref{cap:p1c1}\\
Sistema 1 e Sistema 2 sono due sistemi cerebrali distinti & Superata (utile come distinzione tra processi) & \ref{cap:p1c1}, \ref{cap:p1c2}\\
L'amigdala è «il centro della paura» & Discussa (rete distribuita) & \ref{cap:p1c2}\\
La dissonanza cognitiva spinge a ridurre le contraddizioni & Buona & \ref{cap:p1c22}\\
La saggezza delle folle richiede giudizi indipendenti & Buona & \ref{cap:p1c9}\\
Gli algoritmi di raccomandazione aumentano la polarizzazione di tutti gli utenti & Discussa & \ref{cap:p2c11}\\
I giudici decidono peggio prima della pausa pranzo & Discussa & \ref{cap:p2c28}\\
Strumenti di punteggio come COMPAS producono errori distribuiti in modo diseguale tra gruppi & Buona & \ref{cap:p2c28}\\
Il 2024 è stato l'anno più caldo mai registrato & Solida & \ref{cap:p2c18}\\
Le differenze cerebrali tra i sessi sono medie, con ampia sovrapposizione & Buona & \ref{cap:p2c20}\\
La «minaccia dello stereotipo» peggiora le prestazioni & Discussa & \ref{cap:p2c20}\\
Esistono differenze culturali medie negli stili di attribuzione & Buona & \ref{cap:p2c13}\\
I neuroni specchio spiegano la trasmissione dei pregiudizi & Ipotesi (non supportata) & \ref{cap:p2c12}\\
I bambini apprendono preferenze sociali da segnali non verbali & Buona & \ref{cap:p2c12}\\
L'uso dei social media danneggia in media il benessere degli adolescenti & Discussa & \ref{cap:p2c29}\\
Troppa scelta riduce sempre la probabilità di acquisto & Discussa & \ref{cap:p2c30}\\
Rappresentazioni di parole apprese da testi contengono stereotipi umani & Solida & \ref{cap:p3c11}\\
I modelli linguistici ottimizzati sul gradimento tendono alla compiacenza & Buona & \ref{cap:p3c11}, \ref{cap:p3c35}\\
Il bias di automazione riduce il controllo umano sugli errori delle macchine & Buona & \ref{cap:p3c11}, \ref{cap:p3c27}\\
L'uso dell'IA riduce la competenza di chi la usa (dequalificazione) & Ipotesi con primi dati & \ref{cap:p3c27}\\
L'aiuto dell'IA rende i testi creativi più simili tra loro & Buona (singoli studi) & \ref{cap:p3c27}\\
Un tutor IA senza vincoli può peggiorare l'apprendimento & Buona (singoli studi) & \ref{cap:p4c32}\\
Un allenamento mirato può ridurre alcuni bias & Buona (effetti limitati) & \ref{cap:p4c33}\\
Un simulatore può misurare l'ancoraggio da segnali fisiologici & Ipotesi (scenario) & \ref{cap:p4c33}\\
\bottomrule
\end{longtable}
}

\section*{Riferimenti}

\begingroup
\small
\setlength{\parindent}{-1.5em}
\setlength{\leftskip}{1.5em}
\setlength{\parskip}{0.35em}

Apfelbaum, E. P., Pauker, K., Sommers, S. R., e Ambady, N. (2010). In blind pursuit of racial equality? \textit{Psychological Science}, 21(11), 1587-1592.

Bai, X., Wang, A., Sucholutsky, I., e Griffiths, T. L. (2025). Explicitly unbiased large language models still form biased associations. \textit{Proceedings of the National Academy of Sciences}.

Barrett, L. F. (2017). \textit{How Emotions Are Made: The Secret Life of the Brain}. Houghton Mifflin Harcourt, Boston.

Bastani, H., e altri (2025). Generative AI without guardrails can harm learning: Evidence from high school mathematics. \textit{Proceedings of the National Academy of Sciences}.

Bian, L., Leslie, S.-J., e Cimpian, A. (2017). Gender stereotypes about intellectual ability emerge early and influence children's interests. \textit{Science}, 355(6323), 389-391.

Bigler, R. S., Jones, L. C., e Lobliner, D. B. (1997). Social categorization and the formation of intergroup attitudes in children. \textit{Child Development}, 68(3), 530-543.

Bigler, R. S., e Liben, L. S. (2007). Developmental intergroup theory. \textit{Current Directions in Psychological Science}, 16(3), 162-166.

Bolukbasi, T., Chang, K.-W., Zou, J., Saligrama, V., e Kalai, A. (2016). Man is to computer programmer as woman is to homemaker? Debiasing word embeddings. \textit{Advances in Neural Information Processing Systems}, 29. arXiv:1607.06520.

Brady, W. J., Wills, J. A., Jost, J. T., Tucker, J. A., e Van Bavel, J. J. (2017). Emotion shapes the diffusion of moralized content in social networks. \textit{PNAS}, 114(28), 7313-7318.

Budzyń, K., e altri (2025). Endoscopist deskilling risk after exposure to artificial intelligence in colonoscopy. \textit{The Lancet Gastroenterology \& Hepatology}.

Caliskan, A., Bryson, J. J., e Narayanan, A. (2017). Semantics derived automatically from language corpora contain human-like biases. \textit{Science}, 356(6334), 183-186. DOI: 10.1126/science.aal4230.

Cesario, J., Johnson, D. J., e Eisthen, H. L. (2020). Your brain is not an onion with a tiny reptile inside. \textit{Current Directions in Psychological Science}, 29(3), 255-260.

Chesney, R., e Citron, D. (2019). Deep fakes: A looming challenge for privacy, democracy, and national security. \textit{California Law Review}, 107(6), 1753-1820.

Chouldechova, A. (2017). Fair prediction with disparate impact. \textit{Big Data}, 5(2), 153-163.

Cohen, G. L., e Sherman, D. K. (2014). The psychology of change: Self-affirmation and social psychological intervention. \textit{Annual Review of Psychology}, 65, 333-371.

Coyne, S. M., e altri (2016). Pretty as a princess. \textit{Child Development}, 87(6), 1909-1925.

Danziger, S., Levav, J., e Avnaim-Pesso, L. (2011). Extraneous factors in judicial decisions. \textit{PNAS}, 108(17), 6889-6892.

Doshi, A. R., e Hauser, O. P. (2024). Generative AI enhances individual creativity but reduces the collective diversity of novel content. \textit{Science Advances}, 10(28).

Dratsch, T., e altri (2023). Automation bias in mammography. \textit{Radiology}, 307(4).

Dressel, J., e Farid, H. (2018). The accuracy, fairness, and limits of predicting recidivism. \textit{Science Advances}, 4(1).

Feinstein, J. S., Buzza, C., Hurlemann, R., e altri (2013). Fear and panic in humans with bilateral amygdala damage. \textit{Nature Neuroscience}, 16(3), 270-272.

Festinger, L. (1957). \textit{A Theory of Cognitive Dissonance}. Stanford University Press, Stanford.

Flore, P. C., e Wicherts, J. M. (2015). Does stereotype threat influence performance of girls in stereotyped domains? \textit{Journal of School Psychology}, 53(1), 25-44.

Gonen, H., e Goldberg, Y. (2019). Lipstick on a pig. \textit{Proceedings of NAACL-HLT}, 609-614.

Guess, A. M., e altri (2023). How do social media feed algorithms affect attitudes and behavior in an election campaign? \textit{Science}, 381(6656), 398-404.

Guthrie, C., Rachlinski, J. J., e Wistrich, A. J. (2001). Inside the judicial mind. \textit{Cornell Law Review}, 86, 777-830.

Haselton, M. G., e Nettle, D. (2006). The paranoid optimist. \textit{Personality and Social Psychology Review}, 10(1), 47-66.

Heine, S. J., e Hamamura, T. (2007). In search of East Asian self-enhancement. \textit{Personality and Social Psychology Review}, 11(1), 4-27.

Henrich, J., Heine, S. J., e Norenzayan, A. (2010). The weirdest people in the world? \textit{Behavioral and Brain Sciences}, 33(2-3), 61-83.

Hosseinmardi, H., e altri (2021). Examining the consumption of radical content on YouTube. \textit{PNAS}, 118(32).

International Nuclear Safety Advisory Group (1992). \textit{The Chernobyl Accident: Updating of INSAG-1} (INSAG-7). IAEA, Vienna.

Kahan, D. M., e altri (2012). The polarizing impact of science literacy and numeracy on perceived climate change risks. \textit{Nature Climate Change}, 2, 732-735.

Kahneman, D. (2011). \textit{Thinking, Fast and Slow}. Farrar, Straus and Giroux, New York.

Kahneman, D., e Tversky, A. (1973). On the psychology of prediction. \textit{Psychological Review}, 80(4), 237-251.

Kleinberg, J., Mullainathan, S., e Raghavan, M. (2017). Inherent trade-offs in the fair determination of risk scores. \textit{Proceedings of ITCS}.

Kobak, D., González-Márquez, R., Horvát, E.-Á., e Lause, J. (2025). Delving into LLM-assisted writing in biomedical publications through excess vocabulary. \textit{Science Advances}, 11(27).

LeDoux, J. (1996). \textit{The Emotional Brain}. Simon and Schuster, New York.

Loftus, E. F., e Palmer, J. C. (1974). Reconstruction of automobile destruction. \textit{Journal of Verbal Learning and Verbal Behavior}, 13(5), 585-589.

Lord, C. G., Lepper, M. R., e Preston, E. (1984). Considering the opposite. \textit{Journal of Personality and Social Psychology}, 47(6), 1231-1243.

Lorenz, J., Rauhut, H., Schweitzer, F., e Helbing, D. (2011). How social influence can undermine the wisdom of crowd effect. \textit{PNAS}, 108(22), 9020-9025.

Masuda, T., e Nisbett, R. E. (2001). Attending holistically versus analytically. \textit{Journal of Personality and Social Psychology}, 81(5), 922-934.

Mercier, H., e Sperber, D. (2017). \textit{The Enigma of Reason}. Harvard University Press, Cambridge.

Morewedge, C. K., e altri (2015). Debiasing decisions: Improved decision making with a single training intervention. \textit{Policy Insights from the Behavioral and Brain Sciences}, 2(1), 129-140.

Navajas, J., Niella, T., Garbulsky, G., Bahrami, B., e Sigman, M. (2018). Aggregated knowledge from a small number of debates outperforms the wisdom of large crowds. \textit{Nature Human Behaviour}, 2, 126-132.

Nickerson, R. S. (1998). Confirmation bias: A ubiquitous phenomenon in many guises. \textit{Review of General Psychology}, 2(2), 175-220.

Nyhan, B., e altri (2023). Like-minded sources on Facebook are prevalent but not polarizing. \textit{Nature}, 620, 137-144.

Obermeyer, Z., Powers, B., Vogeli, C., e Mullainathan, S. (2019). Dissecting racial bias in an algorithm used to manage the health of populations. \textit{Science}, 366(6464), 447-453. DOI: 10.1126/science.aax2342.

Orben, A., e Przybylski, A. K. (2019). The association between adolescent well-being and digital technology use. \textit{Nature Human Behaviour}, 3(2), 173-182.

Parasuraman, R., e Manzey, D. H. (2010). Complacency and bias in human use of automation. \textit{Human Factors}, 52(3), 381-410.

Pronin, E., Lin, D. Y., e Ross, L. (2002). The bias blind spot. \textit{Personality and Social Psychology Bulletin}, 28(3), 369-381.

Scheibehenne, B., Greifeneder, R., e Todd, P. M. (2010). Can there ever be too many options? \textit{Journal of Consumer Research}, 37(3), 409-425.

Sharma, M., e altri (2024). Towards understanding sycophancy in language models. \textit{Proceedings of ICLR}.

Skinner, A. L., Meltzoff, A. N., e Olson, K. R. (2017). «Catching» social bias. \textit{Psychological Science}, 28(2), 216-224.

Stavrova, O., e Ehlebracht, D. (2019). The cynical genius illusion. \textit{Personality and Social Psychology Bulletin}, 45(2), 254-269.

Stone, J., Aronson, E., Crain, A. L., Winslow, M. P., e Fried, C. B. (1994). Inducing hypocrisy as a means of encouraging young adults to use condoms. \textit{Personality and Social Psychology Bulletin}, 20(1), 116-128.

Supran, G., Rahmstorf, S., e Oreskes, N. (2023). Assessing ExxonMobil's global warming projections. \textit{Science}, 379(6628).

Talhelm, T., e altri (2014). Large-scale psychological differences within China explained by rice versus wheat agriculture. \textit{Science}, 344(6184), 603-608.

Tomkins, A., Zhang, M., e Heavlin, W. D. (2017). Reviewer bias in single- versus double-blind peer review. \textit{PNAS}, 114(48), 12708-12713.

Tversky, A., e Kahneman, D. (1974). Judgment under uncertainty: Heuristics and biases. \textit{Science}, 185(4157), 1124-1131. DOI: 10.1126/science.185.4157.1124.

Woolley, A. W., e altri (2010). Evidence for a collective intelligence factor in the performance of human groups. \textit{Science}, 330(6004), 686-688.

World Meteorological Organization (2025). \textit{State of the Global Climate 2024}. WMO, Ginevra.

\endgroup
